% Transcription — fonds Alexandre Grothendieck, shelfmark 156-6, batch 4 (pages 61-80).
% Established from the facsimile archives/batches/156-6/batch-04.pdf (a copy of
% 156-6.pdf fetched through the project's relay; no cover sheet in the batch,
% so PDF page k = folder page 60+k), per the skill transcribe-grothendieck. The
% folder is ninety-three pages in five batches; this is the fourth. Batches 1
% and 2 of this folder and batches 1-2 of the first draft (folder 156-4, GF
% IV) were read when this batch was drafted; batches 3 and 5, and the later
% batches of 156-4 (where GF IV's own treatment of refinements presumably
% lies: batch 2 of this folder cites GF IV pp. 53-58), did not yet exist, so
% no page-by-page comparison with the first draft could be made here. The
% notation of batch 2 is kept: 𝔉, 𝓜, ℒ, compatibility \between, disjointness
% \parallel, ≪, interior refinement \overset{\circ}{\ll}, the underlined
% triangle \trianglelefteq, « omb » in roman.
%
% Pass: Opus 5.5 (claude-opus-5-5), 2026-09-26 — first pass, unchecked against the pages by a human.
%
% Pages skipped: none. Pages summarised: none. All twenty leaves are his
% mathematics, in ink, fast drafting.
%
% His pagination 61-80 stands at the head of each leaf and coincides with the
% archivists' (noted on page 61). One date in his hand: « 21. juin »,
% underlined, in the left margin of page 64 against the Proposition — a day
% later than the 18-20/06/1986 of the catalogue title; transcribed as a
% \marginal with a note. No other date on these pages.
%
% Numbered runs, recorded where they stand: the atelier axioms At 9 (p. 61),
% At 9' (p. 62), At 9'' (p. 63), At 10 (stated in the long slanted note on
% page 64; referred to as « At 10, p. 64 » in the margin of page 72), At 11 and
% its weakened form At 11' (pp. 67-68), At 12 (p. 76), At 9_Σ (p. 78), At 11
% bis (p. 80); « Atens 2 » and « Atens* 4 » for ensemblist ateliers (p. 61).
% His formula numbers (1.59) (p. 63), (1.60) (p. 64), (1.61) (p. 66, also
% written and struck on that page as the label of an Inf^≤ formula); he
% cites (1.58) « p. 59 » (pp. 62, 70) and (1.48) (p. 72), both in batch 3.
% Corollaries: Cor. 1-3 on pp. 66-67; Corollaire 2 on p. 69 citing « Cor. 3 de
% la p. 66 »; a second « Cor. 3 » on p. 70 (disjoint sums); Cor. 4 (p. 72,
% after a struck « Cor. 4 » on p. 71); Corollaire 5 (p. 72); « Cor. Main »
% (p. 76, « Main » in English); Corollaires 1-2 of At 9_Σ, called « ancien
% At 9'' » and « ancien At 9' » (p. 78). Internal cross-references: « p. 72 »
% (p. 62 margin), « p. 73 » (p. 77), « cor. 2 p. 69 » (p. 77), « p. 66 »
% (p. 69), « p. 59 » (p. 70), all as plain page numbers.
%
% What the batch is about. It continues the axiomatics of an « atelier »
% (𝔉, ≤, ≪) with its multistrates 𝓜 and lieux ℒ. Pages 61-64: « axiomes des
% lieux » — At 9, the disjointness criterion by lieux (X ∥ Y as soon as all
% x ∈ omb X, y ∈ omb Y are disjoint); Atens 2 (distinct lieux compatible) for
% ensemblist ateliers and a corollary (F ∥ G iff their lieux are pairwise
% disjoint); At 9' (compatibility of X, Y from disjointness of the lieux
% outside omb(X ∩ Y)), his remark (1.59) x ∥ y ⟺ x ≬ y and x ≠ y, and the
% « simplicial » variant At 9'' requiring X ∩ Y ∈ 𝓜; At 10 in the page-64
% slanted note. Pages 64-66 (21 June): Proposition — for F' ≪ F and G ≤ F,
% G' = Inf^≪(F', G) exists, is a sub-figure of F', with (1.60)
% G̃' = {X ∈ F̃' | X ≪ G}; proof via a Lemma using At 4 and At 7; corollaries:
% for F' ≤ F it is F' ∧ G; G ↦ F'.G, Sousfig(F) → Sousfig(F') (1.61) commutes
% with Sup and Inf (via φ*: 𝔓_f(F̃) → 𝔓_f(F̃')); Cor. 3, restriction of a
% refinement to a covering family of sub-figures. Pages 67-70: the « principe
% de recollement des raffinements », axiom At 11 (gluing refinements F'_X of
% the strata X of F), its « noyau » At 11' (X ≬ Y, X' ≪ X, Y' ≪ Y agreeing on
% X ∩ Y ⟹ X' ≬ Y'), the construction F̃' = ∪ F̃'_X, gluing on a compatible
% system (Cor. 2), and disjoint sums: Raff(F) × Raff(G) ≅ Raff(F ⨿ G) for
% F ∥ G, whence (1.58). Pages 71-72: discrete figures (all strata are lieux),
% Cor. 4 characterising discrete refinements of F by parts A ⊂ omb(F), and
% Cor. 5 (a proper sub-figure G of F misses some lieu x ≪ F). Pages 73-76: a
% Proposition — F ≬ G, K = F ∩ G, F' ≪ F, G' ≪ G: F' ≬ G' iff F'|K ≬ G'|K,
% and then F' ∨ G' ≪ F ∨ G — proved through two cube-shaped diagrams of ≤
% and ≪ arrows (p. 73 cancelled, p. 74 fair), reduced to a « Cor. Main » (G ⊴
% F, F'_G ≤ G' ⟹ F' ≬ G'), which he then states as At 12. Pages 77-80: the
% implications among At 9-At 12, the notion of Σ-atelier (Σ a full
% subcategory of ordered sets) and At 9_Σ, which subsumes At 9' and At 9''
% as corollaries; two diagrams of implications (p. 79, the first
% cancelled); At 11 bis (filtering Sups of refinements) and the « reduced »
% systems At 9_Σ, At 10, At 11 bis or At 9, At 11 bis, At 12.
%
% Reading hazards fixed once here. As in batch 2: 𝔉 \mathfrak{F}, 𝓜
% \mathcal{M}, ℒ \mathcal{L}, compatibility (crossed chevrons) \between,
% disjointness ∥ \parallel, the underlined triangle \trianglelefteq. His
% script figure letter on pages 65 and 71-72 (a round X, underlined) is set
% \mathfrak{X}; on page 65 a Z is written over an X and is read Z. His
% « 𝔓f » (closed parts) is \mathfrak{P}_{\mathrm{f}}; the disjoint sum (two
% bars) is \amalg; « Sousfig », « Raff », « Inf », « Sup » are set in roman.
% Drawings: the diagrams of pages 69, 73, 74 (two), 79 (two) are described
% in French notes rather than redrawn; the square of page 64 is set in
% tikz-cd with the ≤, ≪ riding headless lines, as drawn; the small schemes
% of pages 62, 65 and 76 are given as lists of relations with a note. His
% underlining is set \emph. Long slanted notes cover the top half of page 64
% (read only in fragments) and the left margins of pages 62, 64, 67, 78, 80;
% they are given as \marginal where something could be read.
%
% Cancellations: page 69 (part of the proof, with a diagram), the lower half
% of page 73, part of page 74, the end of page 78 and the first diagram of
% page 79 are crossed by long strokes; they are transcribed where legible and
% the cancellation noted once each.
%
% Readings flagged and not settled: the insertion in At 9 (p. 61); the
% opening of page 62's At 9' and most of its margin; the sentences around
% At 9'' (p. 63); the first paragraph of page 64 under the slanted note;
% the paragraph after Cor. 3 on page 67; the top of page 69; the long
% paragraph on page 77 listing the axioms; the NB of page 80. Variances on the
% page and not repaired: « omb X, omb Y » for omb F, omb G (p. 62); « F_X »
% for F'_X (p. 68); « K'' » for K' (p. 73); a repeated « Cor. 3 » (pp. 66, 70).
%
\documentclass[11pt,a4paper]{article}
\input{../preamble/grothendieck}

\folder{156-6}
\batch{4}
\pages{61}{80}
\dating{1986}
\watermark{Édition de démonstration}
\foldertitle{[Chapitre] VI. Analysis situs (deuxième mouture) : notes manuscrites (18-20/06/1986)}

\begin{document}

\page{61}
\note{pagination de sa main 61 à 80 en tête des feuillets, qui coïncide avec celle des archivistes}

mais je ne veux pas l'expliciter en axiome ici ; ou je le vois comme cas particulier d'un \emph{axiome de recollement des raffinements}, beaucoup plus général, dont il sera question plus bas. C'est un axiome qui s'énonce sans référence particulière aux lieux. Mais j'ai envie de \ill{} deux axiomes, impliquant \uncertain{eux} les lieux.

\textbf{At 9} Soient $X$ et $Y \in \mathcal{M}$, telles que $\forall\, x \in \mathrm{omb}(X)$, $y \in \mathrm{omb}(Y)$, $x$ et $y$ \add{\uncertain{soient} \ill{}} disjoints : \struck{\ill{} $x \between y$} $x \parallel y$. Alors $X$ et $Y$ \uncertain{sont} \add{\ill{}} disjoints : $X \parallel Y$.
\marginal{Critère de disjonction par les lieux.}\note{« omb » est souligné dans « $\mathrm{omb}(X)$ »}

D'autre part, pour un atelier ensembliste, (et par exemple \uncertain{quasi-ensembliste} !) on pose

\textbf{Atens 2} Si $x, y \in \mathcal{L}$, $x \neq y$, alors $x \underset{\mathcal{M}}{\between} y$

qui se traduit en termes de \struck{$\mathcal{M}$, $\between$} $\mathcal{M} \subset \mathrm{Fig}(\mathcal{L})$, et $\underset{\mathcal{M}}{\between}$, par

\textbf{Atens* \struck{\ill{}} 4} $\forall\, x, \struck{\ill{}}\, y \in \mathcal{L}$, $x \neq y$, \ldots\ $\lbrace\lbrace x \rbrace\rbrace \underset{\mathcal{M}}{\between} \lbrace\lbrace y \rbrace\rbrace$.
\note{« Atens » (atelier ensembliste) est souligné ; le « 4 » suit un chiffre biffé}

\page{62}
Donc

\textbf{Corollaire} Pour un atelier ensembliste (satisfaisant \struck{\ill{}} aussi At 9 et Atens 2) si $F$, $G$ sont deux figures, pour que $F \parallel G$, il f. et s. que $\forall\, x \in \mathrm{omb}\, X$, $y \in \mathrm{omb}\, Y$, on ait $x \parallel y$. \struck{\ill{}}
\note{il écrit « $\mathrm{omb}\, X$, $\mathrm{omb}\, Y$ » là où l'on attend $\mathrm{omb}\, F$, $\mathrm{omb}\, G$}

Dans le cas des ateliers quelconques, on trouve d'autre part

\textbf{Corollaire} Validité de (1.58), si $\mathfrak{F}$ satisfait At 9.
\note{la formule (1.58) est à la page 59, d'après le renvoi de la page 70}

D'autre part, je voudrais un \uncertain{puissant} \ill{} axiome \uncertain{facultatif}

\textbf{At \struck{\ill{}} 9'} Soient $X, Y \in \mathcal{M}$, \struck{$F$} et $F = X \cap Y$ \add{(\uncertain{dans} $\widetilde{F}$ \ill{}} l'une des strates communes). \struck{Supposons que \ill{} pour tout \ill{} de $\mathcal{M}$}, \add{soit} Supposons que pour tout $x \in \mathrm{omb}(X) \smallsetminus \mathrm{omb}(F)$ et tout $y \in \mathrm{omb}(Y) \smallsetminus \mathrm{omb}(F)$ on ait $x \parallel y$ (donc $x \neq y$) [et que \ill{} indique aussi que \struck{\ill{}}
\[
  \mathrm{omb}(X) \cap \mathrm{omb}(Y) = \mathrm{omb}(F) \,]
\]
Alors $X \between Y$.
\note{sous l'énoncé, un petit schéma : $\mathrm{omb}(X)$ et $\mathrm{omb}(Y)$, reliés par deux traits à $\mathrm{omb}(F)$ au-dessous}

\marginal{NB La réciproque \struck{\ill{}} vraie aussi, par axiome de disjonction plus bas p.~72}

Cet axiome est intuitivement \add{\ill{}}, qu'il exprime \struck{la} \add{\uncertain{\ill{}}} compatibilité des multistrates génériques (qui \ill{} \struck{\ill{} points distincts}

\marginal{NB dans $\mathfrak{F}$, les Inf \struck{\ill{}} existent toujours, \ill{} pour $F \cap G$ \ill{} $F, G$ compatibles, il y correspond \ill{} l'intersection \struck{\ill{}} \ill{} de $\mathcal{M}$ \ill{} partie $F$. L'axiome ci-contre \ill{} automatique \ill{} deux figures \ill{} quelconques \ill{}}\note{note oblique dans la marge gauche, lue par fragments}

\page{63}
ingrédients principaux pour décrire l'atelier, i.e.\ $(\mathcal{M}, \leq, \ll, \between)$) en termes ensemblistes, au moyen de la seule relation $\between$ dans $\mathcal{L}$ (puisque \struck{\ill{}} il est clair que pour deux \emph{lieux},
\begin{equation}
  x \parallel y \Longleftrightarrow x \between y \text{ et } x \neq y \tag{1.59}
\end{equation}
\struck{Le seul} Cela mettrait donc un poids particulier sur cette relation $\underset{\mathcal{L}}{\between}$ (ou ce qui revient au même $\underset{\mathcal{L}}{\parallel}$) \struck{dans} \emph{sur} $\mathcal{L}$. Cet \struck{\ill{}}axiome est bien raisonnable, \uncertain{compris} dans des cas non « ensemblistes », mais la seule \uncertain{chose}, c'est qu'il n'est \emph{pas} satisfait dans les cas \uncertain{simpliciaux} les plus importants, cf exemples \uncertain{\ill{}}\ldots. Mais dans ce cas, on peut \uncertain{substituer} une \struck{propriété} propriété qui fera le \uncertain{même} usage

\textbf{At \struck{\ill{}} 9''} Soient $X, Y \in \mathcal{M}$ et $F$ comme dans At 9' 1. Alors
\[
  X \between Y \Longleftrightarrow
  \begin{cases}
    \forall\, x \in \mathrm{omb}\, X \smallsetminus \mathrm{omb}\, F \text{ et } y \in \mathrm{omb}\, Y \smallsetminus \mathrm{omb}\, F, \text{ on a } x \parallel y \\
    F \text{ est « élémentaire » (i.e. } F \in \mathcal{M} \text{)}
  \end{cases}
\]
\note{sous « élémentaire », il ajoute « i.e. c'est une multistrate commune » ; au-dessus, entouré, « \uncertain{réciproquement} » ; la numérotation « At 9' 1 » est sur la page}

\marginal{Variante « \uncertain{simplicial} » de At 9'}

\marginal{NB On peut donner variantes avec toutes autres conditions sur l'ens. ordonné $\widetilde{F}$.}

\page{64}
At 9 est impliqué par chacun des Axiomes At 9', At 9''. (Et il \ill{} \ill{} \ill{} \ill{} \ill{} coïncide \ill{} \ill{}\ldots)
\note{la moitié supérieure de la page est couverte par une longue note oblique, écrite par-dessus le texte, qui en rend la fin illisible}

\marginal{(*) Introduction At 9, At 10, pour \ill{} ces axiomes At 8, At 10 (qui sont \ill{} \ill{} \ill{} \struck{d'existence}, \ill{} : l'existence \ill{} des lieux et des \ill{} \ill{} \ill{}. \struck{Donnons ici deux \ill{}} Axiome de disjonction\ldots\ \textbf{At 10} Soient $X, Y \in \mathcal{M}$\ldots\ \struck{$X \between Y$, $X \neq Y$, alors \ill{} $\mathrm{omb}(F)$ \ill{} disjoints}\ldots\ $x \in \mathrm{omb}(X)^{\circ}$ et $y \in \mathrm{omb}(Y)^{\circ}$ \ldots\ NB \struck{$Y \neq X$ \ldots\ $X \leq Y$} $X \cap Y \neq X, Y$ résulterait aussi de l'axiome du recollement At 11 plus bas --- dans le cas \emph{général} \ill{} $Y \subset X$. Les axiomes At 8, At 9, At 10 \ill{} « axiomes des lieux ». \ldots\ \textbf{At 10} $X, Y \in \mathcal{M}$, $X \neq Y$ \ldots\ $x \parallel y$ \ldots\ (*) \ldots\ cf.\ p.~72}\note{note écrite en oblique en travers du haut de la page, marquée d'un astérisque cerclé ; lue par fragments seulement. C'est là qu'est énoncé l'axiome At 10 (« axiome de disjonction ») auquel renvoie la marge de la page 72 (« At 10, p.~64 »)}

\marginal{21. juin}\note{date de sa main, soulignée, dans la marge gauche en face de la Proposition}

\textbf{Proposition} Soient $\mathfrak{F}$ un atelier At 1 --- \struck{\ill{}} At 7 \uncertain{Soit} $F \in \mathfrak{F}$, $F' \ll F$, $G \leq F$. Alors
\[
  G' = \mathrm{Inf}^{\ll}(F', G)
\]
existe, et $G'$ est une sous-figure de $F'$ : \struck{\ill{}}

\begin{tikzcd}
F' \arrow[r, no head, "\ll" description] & F \\
G' \arrow[u, no head, "\leq" description] \arrow[r, no head, "\ll" description] & G \arrow[u, no head, "\leq" description]
\end{tikzcd}

\begin{equation}
  \widetilde{G'} = \lbrace X \in \widetilde{F'} \mid X \ll G \rbrace \tag{1.60}
\end{equation}
\struck{\ill{}}.
\note{le numéro (1.60) est écrit dans la marge gauche, en face de la formule}

\textbf{Démonstration} La formule (1.60) \struck{résulte} est tautologique par la définition des $\mathrm{Inf}^{\ll}$ \ill{} le fait que $\leq$ implique $\ll$. Mais sans admettre l'existence des Inf, la formule (1.60) définit une partie fermée de $\widetilde{F'}$, donc une sous-figure $G'$ de $F'$, et on a évidemment $G' \ll F$ \uncertain{par}

\marginal{Peut-être remplacer \ill{} \ill{} \ill{} \ill{} (1.60) \ill{} \ill{} \ill{}}\note{notes obliques dans la marge gauche du bas de la page, lues par fragments}

\page{65}
At 4. Prouvons que c'est un Inf pour $\ll$, i.e.\ que pour toute figure $\mathfrak{X}$, on a
\[
  \mathfrak{X} \ll F',\ \mathfrak{X} \ll G \Longrightarrow \mathfrak{X} \ll G'
\]
Par At 4, on est ramené au cas où $\mathfrak{X} = Z \in \mathcal{M}$. Comme $Z \ll F'$, soit $X'$ le plus petit élément de $\widetilde{F'}{}^{Z} = \lbrace X' \in \widetilde{F'} \mid Z \ll X' \rbrace$. Il suffit de prouver que $X' \ll G'$, i.e.\ $X' \ll G$. Donc on est ramené au
\note{$\mathfrak{X}$ rend sa lettre de ronde soulignée ; le $Z$ est écrit par-dessus un $X$, ici et plus bas}

\textbf{Lemme} Soit\note{schéma : $F' \ll F$ sur la première ligne ; sous $F'$, $X' \leq F'$ ; sous $F$, $G \leq F$ ; sous $X'$, $Z \overset{\circ}{\ll} X'$ (signe pointé d'un petit cercle), et $Z \ll G$ sur la dernière ligne}
\[
  F' \ll F, \quad X' \leq F', \quad G \leq F, \quad Z \overset{\circ}{\ll} X', \quad Z \ll G,
\]
alors \struck{\ill{}} $X' \ll G$.

En effet, soit $X$ l'unique élément de $\widetilde{F}$ tel que $X' \overset{\circ}{\ll} X$. Alors par At 7 on a que
\[
  Z \overset{\circ}{\ll} X
\]
\note{le $Z$ est écrit par-dessus une lettre biffée}
donc $X$ est le plus petit élément de $\widetilde{F}{}^{Z}$. Comme $Z \ll G$, $\exists\, Y \in \widetilde{G}$ avec $Z \ll Y$, \add{(par At 4)} donc comme $\widetilde{G} \subset \widetilde{F}$ d'où $Y \in \widetilde{F}{}^{Z}$, on a par définition de $X$, $X \leq Y$, \struck{d'où $X \in \widetilde{G}$}, donc \struck{$Z \ll X \leq$} comme $X' \ll X \leq Y \leq G$ d'où $X' \ll G$, qed.

\page{66}
\textbf{Corollaire 1} \struck{L'application \ill{}} Si $F'$ est aussi une sous-figure i.e.\ $F' \leq F$, alors $G'$ est une \struck{(1.61) $G' = \mathrm{Inf}^{\leq}(F', G)$} sous-figure de $G$, \ill{} de $F$, et c'est $F' \wedge G$ (sous-figure intersection).

Formule, en remarquant \uncertain{les rôles} de $F'$ et $G$ --- et le fait que $G' = F' \wedge G$ se voit sur (1.60).

\textbf{Corollaire 2} Si $F' \ll F$, \struck{L'application} $G \longmapsto G' \overset{\mathrm{df}}{=} F'.G$
\begin{equation}
  \mathrm{Sousfig}(F) \longrightarrow \mathrm{Sousfig}(F') \tag{1.61}
\end{equation}
commute aux Sup et aux Inf.
\marginal{Notation $F'.G$}

\struck{En} En effet, en termes des applications canoniques croissantes
\[
  \varphi : \widetilde{F'} \longrightarrow \widetilde{F}
\]
d'ens.\ ordonnés, cette application s'identifie à l'application correspondante des \uncertain{sous-ensembles}
\[
  \varphi^{*} : \mathfrak{P}_{\mathrm{f}}(\widetilde{F}) \longrightarrow \mathfrak{P}_{\mathrm{f}}(\widetilde{F'}), \qquad A \longmapsto \varphi^{-1}(A) \,)
\]
qui a les propriétés en question.
\note{$\mathfrak{P}_{\mathrm{f}}$ rend son « $\mathfrak{P}f$ » : l'ensemble des parties fermées}

\textbf{Corollaire 3} Soit $F' \ll F$, et soit $(F_i)$ une famille de sous-figures de $F$ telles que $F = \operatorname{Sup} F_i$. Soit $F_{ij} = F_i \cap F_j$. Considérons les
\marginal{Notation $\vee F_i$}

\page{67}
sous-figures $F'_i$, $F'_{ij}$ de $F'$ induites par les $F_i$, $F_{ij}$. Alors \struck{$F' = \operatorname{Sup}^{\leq} F'_i$, $F'_i \cap F'_j = F'_{ij}$}.
\begin{itemize}
  \item[(i)] Chaque $F'_i$ est un raffinement de $F_i$
  \item[(ii)] Le raffinement induit par $F'_i$ sur $F_{ij}$ \add{($= F_i \cap \uncertain{F_j}$)} (est $F'_{ij}$) est égal à celui induit par $F'_j$ sur $F_{ij}$)
  \item[(iii)] $F' = \operatorname{Sup} F'_i$, et c'est un raffinement de $F$.
\end{itemize}
\note{un trait vertical, terminé en flèche, relie (ii) à (iii) dans la marge}

\ill{} donne un \emph{principe de recollement} des raffinements de $F$ : un tel raffinement est connu (\uncertain{par} (iii)) quand on connaît une « famille compatible » de raffinements des $F_i$. La question est \uncertain{si} \ill{} donnée celle-ci arbitraire\ldots\ \ill{} \ill{} \ill{}. C'est ce qui va être assuré par l'axiome suivant

\marginal{NB Si pour tout $i, j$ \ill{} $F_{ij}$ \ill{} \ill{} ($F'_{ij}$ \ill{} $F'_i|F_{ij}$) \ill{} \ill{} (ii) signifie $F'_i|F_{ij} = F'_j|F_{ij}$}\note{note oblique dans la marge gauche, lue par fragments}

\textbf{At 11} (\emph{Axiome de recollement des raffinements}) Soit $F \in \mathfrak{F}$, et pour tout $X \in \widetilde{F}$, soit $F'_X \ll X$ \struck{\ill{}} un raffinement de $X$. On suppose que pour $X, Y \in \widetilde{F}$, $Y < X$ \struck{$X \neq Y$, \ill{} $X \angle Y$}, on ait \struck{$F' | X \cap Y =$}
\[
  F'_Y = F'_X | Y
\]
Alors $\exists\, F' \ll F$, tel que pour tout $X \in \widetilde{F}$, $F'_X = F' | X$.
\note{le « 11 » de At 11 est écrit par-dessus un autre chiffre}

\page{68}
\textbf{NB} Par le corollaire précédent, $F'$ est unique. On l'appelle le raffinement de $F$ \emph{obtenu par recollement} à partir des raffinements $F'_X$ de ses multistrates.
\note{il écrit « $F_X$ »}

\textbf{NB} \struck{Cet axiome} \add{Le \struck{\ill{}} contenu essentiel de} \struck{\ill{}} cet axiome, c'est que les $F'_X$ \add{étant donnés} soient compatibles entre eux, \struck{\ill{}} de façon à pouvoir former
\[
  \operatorname*{Sup}_{X \in \widetilde{F}} F'_X = F'
\]
(du moins, on pourrait le faire, indépendamment de At 11, si on suppose $\widetilde{F}$ fini\ldots). La difficulté est de comparer $F'_X$ et $F'_Y$ quand on n'a pas $Y \leq X$ ou $X \leq Y$.

\uncertain{Donc} le \uncertain{noyau} de At 11 est ceci :

\textbf{Cor.} (\emph{At 11'}) Soient $X, Y \in \mathcal{M}$ avec $X \between Y$, soit $X' \ll X$, $Y' \ll Y$, \struck{$F = X \cup Y$}, \struck{$Z$} $T = X \cap Y$. Supposons que
\[
  X' | T = Y' | T
\]
(ou, ce qui revient au même, pour toute $Z \in \widetilde{X} \cap \widetilde{Y}$, $X' | Z = Y' | Z$). Alors $X' \between Y'$ !
\note{« Cor. » est surchargé et la mention « (At 11') » est ajoutée au-dessus ; les $X'$, $Y'$ sont écrits par-dessus d'autres lettres}

Ceci est une variante \uncertain{précisée} affaiblie de At 11. Elle implique At 11 \ill{} pour $\widetilde{F}$ fini --- plus généralement, au moins pour At 11 que les $F'_X$ sont \uncertain{compatibles} entre eux, et si on \ill{} $\operatorname{Sup} F'_X$ existe

\page{69}
(ce qui est automatique si \ill{} \ill{} $\widetilde{F}$ fini) donc je dis que les $F'_X$ \uncertain{sont} \ill{} $\operatorname{Sup} F'$ \uncertain{répond} à la question. En \ill{}, \ill{} \ill{}.
\[
  \widetilde{F'} = \bigcup_{X} \widetilde{F'_X}
\]
et on \ill{} $F' \ll F$ par At 4. Prouvons \add{$\forall\, X \in \widetilde{F}$} que $F' | X = F'_X$, \ill{} \ill{} que
\[
  F'_X = \operatorname*{Sup}_{Z \in \widetilde{F}} (F'_Z | X)
\]
Et par la condition \ill{} sur les $F'_Z$ \ill{} \ill{}
\note{la suite est un passage encadré d'une longue courbe et en partie barré de traits obliques : un schéma $F' \ll F$, $F'_Z \ll Z$, $X$, avec $X \cap Z$ au-dessous, et deux lignes de relations $\operatorname{Sup}_Z F'_Z = F' \ll F$, $F'_Z | X \leq F'_X | X \ll X$ ; à droite, barré, « \ill{} raffinements $F'_Z$ de $F$ \ill{} sous-figures $X$ de $F$ \ill{} les supports des $F'_Z$ \ill{} »}

On regarde \ill{} ensembliste
\[
  \widetilde{F'} = \bigcup \widetilde{F'_X} \longrightarrow \widetilde{F} \supset \widetilde{X}
\]
$\widetilde{F}$ est réunion des $\widetilde{X}$, et ici $\widetilde{F'}$ a été obtenu par recollement \struck{des} \add{des systèmes des} $\widetilde{F'_X}$ des \ill{} admissibles des $\widetilde{X}$ \uncertain{recouvrant} la famille admissible $\widetilde{F}$\ldots

\textbf{Corollaire 2} Sous les conditions du Cor.\ 3 de la page 66, si on se donne un \emph{système compatible} ($F'_i$ des raffinements des $F_i$ (définition en termes \ill{})), alors

\page{70}
il existe un raffinement $F'$ de $F$ (nécessairement unique) qui les induit.

Dans le cas où la famille des indices $I$ est \emph{finie}, ceci résulte déjà du seul axiome affaibli \emph{At 11'}.

En particulier

\textbf{Cor.\ 3} Soient $F$, $G$ deux figures avec $F \parallel G$, et considérons la figure \add{$H = F \amalg G$} somme \uncertain{disjointe}. Alors
\begin{itemize}
  \item[(i)] Si $F' \ll F$, $G' \ll G$, on a $F' \parallel G'$, et
  \[
    H' \overset{\mathrm{df}}{=} F' \amalg G' \ll H = F \amalg G .
  \]
  \item[(ii)] \struck{On trouve} On trouve ainsi une bijection
  \[
    \mathrm{Raff}(F) \times \mathrm{Raff}(G) \xrightarrow{\ \sim\ } \mathrm{Raff}(F \amalg G), \qquad (F', G') \longmapsto H' = F' \amalg G'
  \]
  la bijection en sens inverse est donnée par
  \[
    H' \longmapsto (H'|F, H'|G) .
  \]
\end{itemize}
\note{« Cor.\ 3 » répète le numéro du corollaire 3 de la page 66 ; la somme disjointe est notée par deux barres (ici $\amalg$) ; « (ii) » est ajouté devant « On trouve »}

Se généralise au cas d'une famille $F$, \emph{somme} disjointe de sous-familles $F_i$.

On obtient en particulier la relation (1.58) p.~59
\[
  F \parallel G, \quad F' \ll F,\ G' \ll G \Longrightarrow F' \parallel G' !
\]

\page{71}
\struck{Cor.\ 4}

\textbf{Définition} Une figure \struck{\ill{}} $\mathfrak{X}$ est dite \emph{discrète} si toutes ses \struck{multi}strates sont des lieux, i.e.\ $\widetilde{\mathfrak{X}} \subset \mathcal{L}$.

\struck{\ill{}} Ces telles figures sont connues \marginal{\struck{Coroll.}} par la connaissance de la partie $\widetilde{\mathfrak{X}} = A$ de $\mathcal{L}$. Celle-ci est formée de lieux deux à deux disjoints, et pour qu'une partie \emph{finie} de $\mathcal{L}$ « soit » une figure, il f.\ et s.\ que ses lieux soient deux à deux \struck{\ill{} disjoints} disjoints. \struck{\ill{}} Notons que la relation $\leq$ ou $\ll$ sur $\widetilde{\mathfrak{X}} = A$ est la relation discrète (i.e.\ ainsi sur $\mathcal{L}$ !) donc les parties fermées de $A$ sont les parties quelconques, donc les sous-figures correspondent aux parties quelconques.

Plus généralement, si $F$ est une figure, on s'intéresse aux \uncertain{sous-figures} discrètes. Elles sont décrites par certaines \struck{\ill{}} parties $A$ de $\mathcal{L}$.

\page{72}
\textbf{Cor 4} \add{Soit $F$ une figure,} Pour qu'une partie $A$ de $\mathcal{L}$ définisse une figure discrète $\mathfrak{X}$ \struck{et} qui soit $\mathfrak{X} \ll F$, il f.\ et s.\ que
\begin{itemize}
  \item[(i)] $\mathfrak{X} \subset \mathrm{omb}(F)$
  \item[(ii)] $\forall\, X \in \widetilde{F}$, $\mathfrak{X} \cap \mathrm{omb}(X)$ « soit une figure » (p.\,ex.\ soit fini et formé d'éléments deux à deux disjoints).
\end{itemize}

Ceci indique notamment que \ill{} \ill{}, dans l'énoncé de At 9' (\ill{} $X \between Y$), la condition suffisante pour que \ill{} condition \ill{} en termes de lieux. \ill{} \ill{} de disjonction \emph{nécessaire} [mais est visiblement \uncertain{équivalente} à At 10]

\textbf{Corollaire 5} Soient $F$ une figure, $G$ une sous-figure $\neq F$. Alors $\exists\, x \in \mathcal{L}$ tel que $x \ll F$ (i.e.\ $x \in \mathrm{omb}(F)$) et $x \parallel G$.
\marginal{At 10, p.~64.}

Soit $X \in \widetilde{F} \smallsetminus \widetilde{G}$, et $x \in \mathrm{omb}(X)^{\circ}$ (1.48), i.e.\ $x \in \mathcal{L}$, $x \overset{\circ}{\ll} X$. Je dis qu'il fait l'affaire. Pour voir qu'il est disjoint de $G$, on peut \uncertain{utiliser} \uncertain{soit} At 9 (critère de disjonction par les lieux) et At 10 (critère de disjonction pour les lieux), soit se référer à \struck{\ill{}} l'axiome des recollements, \add{construisant} $G \amalg \ill{}$ comme raffinement de $F$\ldots
\note{« (1.48) » renvoie à une formule d'un lot précédent ; le Corollaire 5 est marqué d'un trait vertical dans la marge}

\page{73}
J'ai envie de renforcer At 11', dans la direction suivante :

\textbf{Proposition (?)} Soient $F, G \in \mathfrak{F}$, $F \between G$, et soit $K = F \cap G$. Soient $F' \ll F$, $G' \ll G$. Pour que $F' \between G'$, il f.\ et s.\ que $F' | K \between G' | K$, et alors $F' \vee G' \ll F \vee G$.

Supposons que $F'$ et $G'$ soient compatibles, \struck{soit $K' = F' \wedge G'$} i.e.\ que $H' = F' \vee G'$ existe. Je dis que $H' \ll F \vee G \overset{\mathrm{df}}{=} H$ : cela résulte aussitôt de la condition At 4. D'autre part, $F'|K$ étant une sous-figure de $H'$, et de même pour $G'|K$, donc \struck{soit $K' = F' \wedge G'$}
\[
  K' \leq F' \ll F \quad \text{et} \quad K' \leq G' \ll G, \quad \text{donc} \quad K'' \ll F, G,
\]
donc (\ldots) $K' \ll K = F \cap G = F.G$. Donc \ill{}. Ils sont compatibles. \struck{on a le diagramme suivant} Inversement, supposons $F'|K$ et \struck{compatibles et} $G'|K$ compatibles
\note{à partir d'ici et jusqu'au bas de la page, le texte est barré de longs traits obliques ; il est transcrit tel quel. « $K''$ » est sur la page}

\note{schéma (barré) : un hexagone $F, H, G$ en haut, $F', K', G'$ en bas, avec $K$ et $H'$ au centre ; flèches à crochet pour les $\leq$ ($F \to H$, $G \to H$, $K \to F$, $K \to G$, $K' \to F'$, $K' \to G'$, $F' \to H'$, $G' \to H'$) et flèches verticales pour les $\ll$ ($F' \to F$, $G' \to G$, $H' \to H$, $K' \to K$). Il est repris au propre page 74}

(les deux flèches horizontales \ill{} \ill{}) --- les 8 flèches horizontales sont des $\leq$, et les grandes flèches verticales des $\ll$ \ill{} plus.

On a
\[
  K' = \mathrm{Inf}^{\ll}(F', G') = \mathrm{Inf}^{\ll}(F', G', K) = \mathrm{Inf}^{\ll}\bigl(\underbrace{\mathrm{Inf}(F', K)}_{F'|K}, \underbrace{\mathrm{Inf}(G', K)}_{G'|K}\bigr)
\]
\struck{$F'.G'$} et $F'|K$ étant une sous-figure de $F' \leq H'$, \ill{} \ill{} de $H'$, et de même pour $G'$

\page{74}
On arrive donc au diagramme
\note{schéma : même disposition qu'à la page 73, au propre. En haut $H$ ; en dessous $F$ et $G$, reliés à $H$ par des flèches à crochet ; au centre $K$, relié à $F$ et à $G$ par des flèches à crochet et à $H$ par une flèche verticale, et une flèche de $K$ vers $H'$ ; puis $F'$ et $G'$, reliés à $F$ et $G$ par des flèches verticales ; en bas $K'$, relié à $F'$ et $G'$ par des flèches à crochet et à $K$ par une flèche verticale. Un premier « $F$ » en haut à gauche est biffé}

\struck{où $K' = F' \cap G'$}

où $K' = F' \cap G'$. (NB.\ $K' \ll K$ car $K' \ll F$ et $K' \ll G$, $K = \mathrm{Inf}^{\ll}(F, G)$) \struck{On a donc} Mieux,
\[
  K' = \mathrm{Inf}^{\ll}(F', G') = \mathrm{Inf}^{\ll}(F', G', K) = \mathrm{Inf}^{\ll}\bigl(\mathrm{Inf}^{\ll}(F', K), \mathrm{Inf}^{\ll}(G', K)\bigr)
\]
i.e.
\[
  K' = \mathrm{Inf}(\underbrace{F'|K, G'|K}_{\text{deux sous-figures de } H'}) = (\underbrace{F'|K}_{F'_K}) \cap (\underbrace{G'|K}_{G'_K}) .
\]
\note{sous « deux sous-figures », un mot biffé}

\struck{et je dis que}
\[
  H'|F = F' \vee G'_K, \qquad H'|G = G' \vee F'_K
\]
\[
  H'_F = F' \vee H'_K, \qquad = G' \vee H'_K
\]
\note{les deux premières égalités sont traversées de traits courbes (biffées ?) ; « $H'|F$ » est surmonté de « $H'_F$ » ; la seconde colonne de la dernière ligne est pour $H'_G$}

on arrive :
\note{schéma : un cube tronqué, de haut en bas : $H$ ; $F$ et $G$ ; $K$ et $H'$ au centre ; $H'_F$ et $H'_G$ ; $F'$, $H'_K$, $G'$ ; $F'_K$ et $G'_K$ ; $K'$. Flèches à crochet ($\leq$) : $F \to H$, $G \to H$, $K \to F$, $K \to G$, $F' \to H'_F$, $G' \to H'_G$, $H'_F \to H'$, $H'_G \to H'$, $H'_K \to H'_F$, $H'_K \to H'_G$, $F'_K \to F'$, $F'_K \to H'_K$, $G'_K \to H'_K$, $G'_K \to G'$, $K' \to F'_K$, $K' \to G'_K$ ; flèches verticales ($\ll$) : $H' \to H$, $H'_F \to F$, $H'_G \to G$, $H'_K \to K$}

\marginal{Diagramme où les cinq carrés horizontaux sont \struck{\ill{} des $\leq$} cartésiens et cocartésiens (pour flèches $\xrightarrow{\leq}$) et les \struck{\ill{}} quatre flèches verticales sont des $\xrightarrow{\ll}$ \ill{} \ill{}.}

\page{75}
Inversement, supposons que
\[
  F' \ll F \quad \text{et} \quad G' \ll G
\]
soient tels que $F'_K \between G'_K$ \struck{\ill{}, désignons par $H'$}.
\[
  H'_K = F'_K \vee G'_K \qquad \text{on aura } H'_K \ll K \text{ par At 4}
\]
\[
  K' = \underbrace{F'_K \cap G'_K}
\]
\struck{\ill{}}

Considérons $F'$ et $H'_K$ \struck{deux raffinements compatibles de} raffinements de $F$ l'un, de $K$ l'autre, et tels que $F'_K \overset{\mathrm{df}}{=} F'|K$ soit compatible avec $H'_K$. Si on savait déjà qu'ils sont compatibles, on pourrait construire leur \struck{Sup} Sup $H'_F$ et $H'_F \to F$, et symétriquement $H'_G$ (\struck{Sup} de $G'$ et de $H'_K$). On vérifierait sans mal $H'_F | K = H'_K$, et de même $H'_G | K = H'_K$, donc par l'axiome de recollement, on saurait $H'_F \between H'_G$ et on pourrait construire $H' = H'_F \vee H'_G$. \struck{Dém.} \uncertain{Nous voyons}, pour vérifier la partie non triviale de la proposition, on est ramené au cas particulier où $G \trianglelefteq F$ et où $F'_G \leq G'$.
\note{sous $G$ et $F$ il écrit $G'$ et $F'$ avec le signe $\ll$ tourné : $G' \ll G$, $F' \ll F$}

\page{76}
i.e.\ on est ramené au

\textbf{Cor.\ Main} Soient $G \trianglelefteq F$, $F' \ll F$ un raffinement de $F$, induisant le raffinement $F'_G \leq$ \struck{$\ll$} $F'$ de $G$, et \struck{soit} supposons $F'_G \leq G'$, où $G'$ est un raffinement de $G$. Alors \struck{$G'$} $F'$ et $G'$ sont compatibles (et $F' \cap G' = F'_G$) \struck{\ill{}}.
\note{« Main » en anglais, souligné avec « Cor. »}

Est-ce que ça résulte des autres axiomes ? \struck{Regardons} On peut encore \uncertain{visualiser} plus (c'est standard) de façon \ill{} suivante

\struck{At 12} \textbf{Corollaire} \add{Main} Soient
\note{schéma : $X \supseteq Y$ sur la première ligne ; sous $X$, $X' \overset{\circ}{\ll} X$ ; sous $Y$, $X'_Y \ll Y$ ; à droite, $Y' \overset{\circ}{\ll} Y$ ; sur la ligne du bas $X' \supseteq X'_Y$, avec $X'_Y = X'|Y$}
\[
  X \supseteq Y, \quad X' \overset{\circ}{\ll} X, \quad Y' \overset{\circ}{\ll} Y, \quad X' \supseteq X'_Y = X'|Y .
\]
Si $X'_Y \between Y'$, alors $X' \between Y'$.

C'est un renforcement \add{(\ill{} cas particulier de At 11 ou At 11')} \ill{} \ill{} des At 10, mais qui résulterait aussi de At 9'' $+$ At 10, ou de At 9' $+$ At 10.

On va formuler plutôt, sous forme un peu plus générale, qui \ill{} impliquera \uncertain{aussi} At 11' (donc, en tout, \uncertain{aussi} At 11) :

\textbf{At 12} Soient $X' \ll X$ et $Y' \ll Y$ dans $\mathcal{M}$, avec $X \between Y$. Soit $K = X \cap Y$, d'où
\[
  X'_K \leq X', \quad X'_K \ll K \quad \text{et} \quad Y'_K \leq Y', \quad Y'_K \ll K .
\]
\struck{Supposons que} Alors

\page{77}
si $X'_K \between Y'_K$, alors $X' \between Y'$. En d'autres termes, si $\forall\, X'' \in \widetilde{X'}$ tel que $X'' \ll K$ ($\Leftrightarrow X'' \ll Y$), et tout $Y'' \in \widetilde{Y'}$ tel que $Y'' \ll K$ ($\Leftrightarrow Y'' \ll X$), on a $X'' \between Y''$, alors $X' \between Y'$.

Ceci implique At 11', At 10, et la prop.\ de la p.~73. Donc cela implique At 11 pour \struck{familles finies} $\widetilde{F}$ \emph{fini}, et le cor.\ 2 p.~69 pour \ill{} \emph{finie}. \struck{\ill{} un énoncé qui implique}

\marginal{\ill{} \ill{} \ill{} At 10,}

\struck{\ill{}} Les axiomes At 9 (critère de disjonction \add{par les lieux}) \add{variantes 9', 9'' (critères de compatibilité par les lieux)} At 10 (critère de disjonction \add{des lieux}) At 11 (critère de recollement des \ill{} \add{(variante un peu affaiblie At 11')}) At 12 (critère de compatibilité par raffinements) \ill{} dans le même \ill{} par \ill{} \ill{} d'\uncertain{certains} de figures (figures de \ill{}, \ill{} les \ill{} \ill{}) \uncertain{intimement} \ill{}, reliés entre eux, par des implications diverses et diagrammes. Mais quitte \add{d'abord} à \add{résumer} une formulation commune, je vais \ill{} à At 9' et At 9'', par la notion de $\Sigma$-atelier, où $\Sigma$ est une sous-cat.\ pleine de la cat.\ des ens.\ ordonnés. Un atelier est un $\Sigma$-atelier si
\begin{itemize}
  \item[(i)] $\forall\, \struck{X}\ X \in \mathcal{M}$, $\widetilde{X} \in \Sigma$
  \item[(ii)] $\mathfrak{F}$ satisfait l'axiome suivant :
\end{itemize}
\note{un bloc biffé de hachures précède « Les axiomes »}

\page{78}
\textbf{At $9_{\Sigma}$} (critère de compatibilité par les lieux, pour $\Sigma$-ateliers) Soient $X, Y \in \mathcal{M}$, \struck{tels que} \struck{\ill{}}
\[
  K = X \cap Y \quad \Bigl( = \mathrm{Inf}^{\leq}(X, Y) = \operatorname*{Sup}_{Z \in \widetilde{X} \cap \widetilde{Y}} Z \Bigr),
\]
de sorte que $\widetilde{K} = \widetilde{X} \cap \widetilde{Y}$. Pour que $X \between Y$, \struck{\ill{}} il f.\ et s.\ que
\begin{itemize}
  \item[(i)] $\forall\, x \in \mathrm{omb}(X) \smallsetminus \mathrm{omb}(K)$, $y \in \mathrm{omb}(Y) \smallsetminus \mathrm{omb}(K)$, on a $x \parallel y$ (\uncertain{ou encore} $x \between y$)
  \item[(ii)] $\widetilde{K} \in \Sigma$
\end{itemize}
\note{à la fin du (i), le signe de compatibilité est barré d'un trait : on attend « $x \between y$ et $x \neq y$ »}

Cela implique donc

\textbf{Corollaire 1} (ancien At 9'') Si on a (i) et
\begin{itemize}
  \item[(ii')] $K \in \mathcal{M}$ (i.e.\ $\widetilde{K}$ a un plus grand élément),
\end{itemize}
alors $X \between Y$.

En effet, $K \in \mathcal{M} \Rightarrow \widetilde{K} \in \Sigma$.

\textbf{Corollaire 2} (ancien At 9') Supposons que $\Sigma \supset$ catégorie des ens.\ ordonnés isomorphes à un $\widetilde{F}$, $F \in \mathfrak{F}$ (p.\,ex.\ $\Sigma =$ cat.\ de tous les ens.\ ordonnés) alors $X \between Y$ ssi on a (i).

Tous les ateliers auxquels j'ai songé \struck{\ill{} supposons que \ill{} les ateliers} \add{sont des} $\Sigma$-ateliers, pour $\Sigma$ convenable, \struck{\ill{}} ce qui signifie que c'est un $\Sigma$-atelier pour $\Sigma = \Sigma_{\mathcal{M}} =$ cat.\ des ens.\ ord.\ \ill{} $\widetilde{X}$, $X \in \mathcal{M}$.
Cela signifie aussi que At $9_{\Sigma}$ est valable pour $\Sigma_{\mathcal{M}}$, on en conclut que l'on \ill{} \ill{} la validité du cor.\ 1
\note{les six dernières lignes sont barrées de grandes croix}

\marginal{\uncertain{Justification} : 1°) $\Sigma =$ tous \ill{} \ill{} 2°) $\Sigma = \Sigma_{\mathcal{M}}$ \ill{} cat.\ \ill{} (ii) \ill{} $X \in \mathcal{M}$}\note{note oblique dans la marge gauche, lue par fragments}

\page{79}
Ceci posé, on a :
\[
  \begin{array}{lll}
    \text{At } 9_{\Sigma} \Longrightarrow \text{At } 9 & & \text{At } 11 \Longleftrightarrow \text{At } 11 \text{ pour } F = \operatorname{Sup} F_i \text{ quelc.} \\
    \Downarrow & & \\
    \text{At } 12 \Longrightarrow \text{At } 10 & & \text{At } 11' \Longleftrightarrow \text{At } 11 \text{ pour } F = \operatorname{Sup} F_i \text{ fini} \\
    \text{At } 12 \Longrightarrow \text{At } 11' \text{ pour } \widetilde{F} \text{ fini} & &
  \end{array}
\]
\note{schéma d'implications : doubles flèches de At $9_{\Sigma}$ vers At 9 et vers At 12, de At 12 vers At 10 et vers « At 11' pour $\widetilde{F}$ fini », de At 11 vers At 11', de « At 11 pour $F = \operatorname{Sup} F_i$ quelc. » vers « At 11 pour $F = \operatorname{Sup} F_i$ fini », et de At 11' vers « At 11' pour $\widetilde{F}$ fini ». Une équivalence avec « At 11 » à la fin de la dernière ligne est biffée, ainsi qu'un signe dans la marge gauche}

Dans ces axiomes, le seul qui fait \ill{} la notion d'un \struck{\ill{}} \add{critère} d'existence pour un \struck{\ill{}} Sup d'une famille \emph{infinie} de figures \struck{condition} est At 11 (qui par cette raison n'est pas \ill{} par At 12) --- la seule autre qui \ill{} \ill{} \ill{} un sup pour un $\mathfrak{F}$
\note{suit un premier schéma d'implications (At $9_{\Sigma}$, At 9, At $9_{\Sigma}$ $+$ At 10, At 12, At 10, At 11, At 11', avec les recollements des raffinements dans $F = \operatorname{Sup} F_i$), barré de grands traits obliques ; il est repris au propre au-dessous}

\note{schéma d'implications : de « At $9_{\Sigma}$ $+$ At 10 » partent deux doubles flèches, vers At $9_{\Sigma}$ et vers At 12 ; de At $9_{\Sigma}$ une double flèche vers At 9 (critère de disj.\ par les lieux) ; de At 12 deux doubles flèches, vers At 10 et vers At 11' ; de At 11 une double flèche vers At 11'. Les légendes sont transcrites ci-dessous}
\begin{itemize}
  \item At 11 ($\Longleftrightarrow$ Recollement des raff.\ dans $F = \operatorname{Sup} F_i$) ($\Longleftrightarrow$ At 11 bis $+$ At 11')
  \item At 11' ($\Longleftrightarrow$ Recollement des raffinements dans $F = \operatorname{Sup} F_i$ pour $I$ fini)
  \item At 9 (critère de disj.\ par les lieux)
\end{itemize}

Le seul axiome qui échappe à At $9_{\Sigma}$ $+$ At 10 (équivalent : At $9_{\Sigma}$ $+$ At 12) est At 11, qui implique un Sup infini. On pourrait donc renforcer At 11 par un énoncé \struck{\ill{}} qui soit équivalent, compte tenu de At 11' :

\page{80}
\textbf{At 11 bis} Soit $F$ une figure, $(F'_i)_{i \in I}$ une famille filtrante croissante de raffinements de $F$. Supposons que pour tout $X \in \widetilde{F}$, la famille des $F'_i | X$ \struck{\ill{}} ait un Sup (soit $F_X$). Alors
\[
  F' = \operatorname{Sup} F'_i
\]
existe.

NB Ça résulte aussitôt de At 11, et \ill{} implique At 11 moyennant At 11' (critère de compatibilité). \struck{Mais ce} \ill{} \ill{} \uncertain{sympathique}, c'est \ill{} le \ill{} « condition de \uncertain{locale} finitude ». Ainsi, on \ill{} : \ill{} aux axiomes « réduits » :

\marginal{n'a rien à voir avec les lieux, et devrait se formuler avec les axiomes sans les lieux (critère d'existence des Sup filtrants)}\note{note de sa main dans la marge droite, reliée par une flèche à At 11 bis}
\[
  \text{At } 9_{\Sigma}, \quad \text{At } 10, \quad \text{At } 11 \text{ bis}
\]
le seul petit ennui étant dans la formulation de At $9_{\Sigma}$, un peu \ill{}, i.e.\ le choix d'un $\Sigma$. Si on veut \ill{} At $9_{\Sigma}$, on est réduit à \uncertain{nouveau} : At 9, \struck{\ill{}}, At 11 bis, At 12
\marginal{ici At 12 disparaît}
\marginal{ici At 10 disparaît ; dans les deux systèmes d'axiomes At 11 et At 11' disparaissent, au profit de At 11 bis qui \ill{} \ill{}}

et on perd seulement un critère de compatibilité universel par lieux, At $9_{\Sigma}$, bien sympathique. Je \ill{} la \ill{} \ill{} pour \ill{} à la saturation de l'atelier, \ill{} n'oubliant pas, dans la définition d'un atelier saturé, d'inclure At 9' $=$ At $9_{\overline{\Sigma}}$ (pour $\overline{\Sigma} =$ tous les ens.\ ordonnés).

\end{document}
